\section{Local regularity for the rank-one trace equation}
\label{reg-section}

The estimate needed for the scalar trace equation is stronger than a
general measurable-coefficient nondivergence estimate. Its distinguished
eigendirection is the gradient of the solution. We prove the resulting
gradient estimate first, and then use its Hessian measure bound in the
second scalar equation. All the estimates in this section are local in
dimension three.

\begin{theorem}[Rank-one gradient and Hessian estimates]
\label{reg-rank-one}
Let $U$ be an interior coordinate ball or a smooth boundary half-ball,
with a uniformly controlled $C^2$ metric $g$. In the boundary case use
Gaussian coordinates, so $g_{33}=1$ and $g_{a3}=0$ for $a=1,2$.
Suppose that $z$ is smooth on $U$, up to the face on its original side
in the boundary case, and satisfies
\begin{equation}
 A^{ij}\nabla_i\nabla_j z=G,\qquad
 A^{ij}=g^{ij}-(1-b)e^i e^j,\qquad
 e=\frac{\nabla z}{|\nabla z|},\qquad c_0\le b\le1,
 \label{reg-trace-equation}
\end{equation}
where $c_0>0$. In the boundary case assume that $z$ is constant on
the face. At a zero of $\nabla z$, an arbitrary measurable unit
direction is allowed, provided Equation~\eqref{reg-trace-equation}
holds with that choice. Assume $|\nabla z|\le L$ and $|G|\le Q$.
There are $\alpha\in(0,1)$, a uniform positive radius $r_0$, and $C$,
depending only on $c_0$, the indicated geometry, and separation from
the other faces of the patch, such that on a smaller patch
\begin{equation}
 \|Dz\|_{C^{0,\alpha}}\le C(L+Q),\qquad
 \int_{B_r(x)\cap U}|\Hess z|^2\,dV_g
       \le C(L+Q)^2r^{1+2\alpha}\quad(0<r\le r_0).
 \label{reg-main-estimates}
\end{equation}
The constants $\alpha,r_0,C$ do not depend on a modulus of continuity
of $b$ or $G$.
The same assertion holds for uniformly controlled families of patches.
\end{theorem}

Here and below, smoothness can be replaced by the piecewise smooth
regularity produced by the reflection described in the proof. In
particular, the estimates do not assume a uniform bound for the Hessian.

\subsection{A cutoff estimate and the gradient flux identity}

\begin{lemma}[A uniform exponent above two]
\label{reg-cutoff}
There is $p_0\in(2,3)$ depending only on $c_0$ with the following
property. If $g$ is sufficiently close to the Euclidean metric in
$L^\infty$, with bounded first derivatives, then the operator in
Equation~\eqref{reg-trace-equation}, with its coefficients held fixed,
satisfies the interior estimate
\begin{equation}
 \|v\|_{W^{2,p_0}(B_{1/2})}
 \le C\bigl(\|v\|_{L^{p_0}(B_1)}
                 +\|A^{ij}\nabla_i\nabla_jv\|_{L^{p_0}(B_1)}\bigr).
 \label{reg-cutoff-estimate}
\end{equation}
It also satisfies the analogous estimate with exponent two.
The coefficient direction in this assertion need not be the gradient
direction of $v$.
\end{lemma}

\begin{proof}
For the full Hessian array the Fourier multiplier
$D^2\Delta^{-1}$ has $L^2$ operator norm one, since
$\sum_{i,j}\xi_i^2\xi_j^2/|\xi|^4=1$. Its finite $L^p$ bounds and
interpolation give constants $C_p$ with $C_p\longrightarrow1$ as
$p\downarrow2$. In Euclidean coordinates the Frobenius norm of
$I-A$ is exactly $1-b$. A sufficiently small perturbation of the
metric therefore gives
$|I-A|\le1-c_0/2$. Choose $p_0\in(2,3)$ so that
$C_{p_0}(1-c_0/2)<1$. For a compactly supported $v$,
\[
 \|D^2v\|_{p_0}
 \le C_{p_0}\|A^{ij}D_{ij}v\|_{p_0}
       +C_{p_0}(1-c_0/2)\|D^2v\|_{p_0}.
\]
Absorb the last term. Christoffel symbols add a bounded first-order
term. Applying this inequality to cutoffs on nested balls and using
$\|Dv\|_p\le\varepsilon\|D^2v\|_p+C_\varepsilon\|v\|_p$
gives Equation~\eqref{reg-cutoff-estimate}; the usual nested-radius
absorption removes the outer Hessian term. The same argument at
$p=2$ proves the last assertion. All coefficient comparisons are
pointwise, so no derivatives of $b$ or of its axis have been used.
\end{proof}

Put $P=\nabla z$, $H=\Hess z$, and $w=|P|^2/2$. A cancellation
specific to the rank-one matrix gives
\begin{equation}
 A\nabla w-GP=(H-g\Delta z)P,
 \qquad
 \divg\bigl(A\nabla w-GP\bigr)
   =|H|^2-(\Delta z)^2+\Ric(P,P).
 \label{reg-flux-identity}
\end{equation}
Indeed, $\nabla w=HP$ and
$G=\Delta z-(1-b)H(e,e)$; thus the terms containing $1-b$
cancel before differentiating. At $P=0$ both cancelled vectors vanish.
The second identity follows by differentiating the remaining tensor
expression and commuting the covariant derivatives of a scalar.
In particular it is meaningful even when $b$ is measurable.
Since
\[
 \Delta z=G+(1-b)H(e,e),
\]
choose $a>0$ so that $(1+a)(1-c_0)^2<1$. Young's Inequality gives
\begin{equation}
 |H|^2-(\Delta z)^2\ge c|H|^2-CG^2.
 \label{reg-hessian-coercivity}
\end{equation}
For $c_0=1$ this follows directly from $\Delta z=G$.

We specify how Equations~\eqref{reg-flux-identity} and
\eqref{reg-hessian-coercivity} are used at a boundary. Subtract the
constant boundary value, extend $z$ oddly across $x_3=0$, extend
$g_{ab}$ and $b$ evenly, and extend $G$ oddly. Keep $g_{33}=1$ and
$g_{a3}=0$. The extension of $z$ is $C^1$: its tangential derivatives
vanish on the face and its normal derivative extends evenly. Hence
its distributional Hessian has no surface atom. The doubled metric
is Lipschitz and smooth on each side; we do not assert that its Ricci
tensor is a bounded function across the joining plane.

Instead apply Equation~\eqref{reg-flux-identity} separately on the
two sides. If $J=(H-g\Delta z)P$, its one-sided normal component
on the face is
\begin{equation}
 J_\nu=-z_\nu\tr_{T\partial U}H
       =\pm H_{\partial U}z_\nu^2.
 \label{reg-boundary-flux}
\end{equation}
The sign depends on the convention for the second fundamental form;
only the bound $|J_\nu|\le |H_{\partial U}|\,|P|^2$ is used.
Consequently the jump of the flux is a plane density bounded by
$2|H_{\partial U}|\,|P|^2$. A bounded density
$q(x')\delta_{\{x_3=d\}}$ is the divergence of the bounded vector
$q(x')\mathbf 1_{\{x_3>d\}}\partial_3$. No tangential derivative
of $q$ is required. Multiplication by $\sqrt{\det g}$ converts the
identity to Euclidean divergence form.

At any chosen center, a constant linear coordinate change makes
$g(0)=I$. In Gaussian charts it can be taken block diagonal, so a
reflected plane remains a plane, possibly displaced from the center.
After a spatial dilation of size $r$, the deviations of $g$ and its
first derivatives from the flat metric are $O(r)$, the ordinary Ricci
terms on each side are $O(r^2)$, and the reflected second fundamental
form is $O(r)$. Normalize the dependent variable so $|\nabla z|\le1$.
If these quantities and $\|G\|_\infty$ are at most $\varepsilon$,
then $s=1-|\nabla z|^2\ge0$ satisfies
\begin{equation}
 -\partial_i(\mathcal A^{ij}\partial_js)
 \ge c|\Hess z|^2-C\varepsilon+\divg E_\varepsilon,
 \qquad
 \|E_\varepsilon\|_\infty\le C\varepsilon,
 \qquad \mathcal A^{ij}=\sqrt{\det g}\,A^{ij}.
 \label{reg-supersolution}
\end{equation}
The divergence here is the coordinate divergence. The vector
$E_\varepsilon$ contains $2\sqrt{\det g}\,GP$ and the bounded
vector representing the plane density. The bulk Ricci and $G^2$
terms have been included in $C\varepsilon$. This proves
Equation~\eqref{reg-supersolution} both in an interior ball and in
a reflected ball, with uniform ellipticity of $\mathcal A$.

\subsection{The fixed-axis equation and improvement of flatness}

\begin{lemma}[Regularity with a fixed axis]
\label{reg-fixed-axis}
Let $a$ be a fixed Euclidean unit vector and let $c_0\le b(x)\le1$
be measurable. If $Y\in W^{2,2}(B_{3/4})$ satisfies
\begin{equation}
 \Delta_{a^\perp}Y+bY_{aa}=0,
 \label{reg-fixed-equation}
\end{equation}
then for some $\alpha_1\in(0,1)$ depending only on $c_0$,
\begin{equation}
 \|Y\|_{C^{1,\alpha_1}(B_{1/4})}
       \le C\|Y\|_{W^{2,2}(B_{3/4})}.
 \label{reg-fixed-axis-estimate}
\end{equation}
\end{lemma}

\begin{proof}
Rotate so that $a=\partial_3$ and put $v=Y_3$. Differentiating
Equation~\eqref{reg-fixed-equation} in the sense of distributions
gives
\[
 \partial_1^2v+\partial_2^2v+\partial_3(b\partial_3v)=0.
\]
This is a scalar uniformly elliptic divergence equation, even when
$b$ depends on all three variables. The scalar De Giorgi--Nash
estimate gives a $C^{0,\alpha_1}$ bound for $v$ on a smaller ball,
controlled by $\|v\|_2$. Caccioppoli's Inequality, applied after
subtracting $v(x)$, then gives for balls contained in that smaller
region
\begin{equation}
 \int_{B_r(x)}|Dv|^2\le Cr^{1+2\alpha_1}
                \|Y\|_{W^{2,2}(B_{3/4})}^2.
 \label{reg-fixed-morrey}
\end{equation}
Now $\Delta Y=F=(1-b)v_3$. By Cauchy--Schwarz,
\begin{equation}
 \int_{B_r(x)}|F|\le Cr^{2+\alpha_1}
                \|Y\|_{W^{2,2}(B_{3/4})}.
 \label{reg-fixed-source}
\end{equation}
Take a cutoff equal to one on $B_{1/3}$, supported where these
estimates hold, and let $V$ be the Newton potential of the cut-off
$F$. Its gradient kernel is bounded by $C|x|^{-2}$ and the derivative
of that kernel by $C|x|^{-3}$. If $d=|x-y|$, the contribution to
$|DV(x)-DV(y)|$ from distances at most $2d$ is bounded, by dyadic
annuli and Equation~\eqref{reg-fixed-source}, by
$C\sum_{2^{-j}\le 2d}(2^{-j})^{\alpha_1}\le Cd^{\alpha_1}$.
At larger distances the bound is
$Cd\sum_{2^{-j}>d}(2^{-j})^{\alpha_1-1}\le Cd^{\alpha_1}$.
The same bounds hold for the cut-off source, with an additional
harmless large-scale term. Thus $DV$ is $C^{0,\alpha_1}$.
The Newton kernel is locally in $L^2$ in dimension three, so
$\|V\|_\infty$ is controlled by $\|F\|_2$. Finally $Y-V$ is
harmonic on $B_{1/3}$; its ordinary interior derivative estimates
complete Equation~\eqref{reg-fixed-axis-estimate}.
\end{proof}

\begin{lemma}[Drop or affine approximation]
\label{reg-dichotomy}
Fix $r_*\in(0,1/32]$. For every $h_*>0$ there are
$k_*\in(r_*,1)$ and $\varepsilon_*>0$ with the following property.
For a normalized solution on $B_1$ in the class of
Equation~\eqref{reg-supersolution}, with errors at most
$\varepsilon_*$, either
\begin{equation}
 \sup_{B_{r_*}}|\nabla z|\le k_* ,
 \label{reg-drop}
\end{equation}
or there is an affine $L$ such that
\begin{equation}
 \sup_{B_{r_*}}|z-L|\le h_*r_* ,
 \qquad \tfrac12\le |DL|\le2.
 \label{reg-affine-alternative}
\end{equation}
If the center lies on the reflecting plane and $z$ has zero boundary
value, $L$ can be chosen odd about that plane.
\end{lemma}

\begin{proof}
Otherwise take errors $\varepsilon_j\to0$ and $k_j\uparrow1$
with neither alternative, and subtract $z_j(0)$ in interior balls.
Lemma~\ref{reg-cutoff}, together with the gradient bound, bounds
$z_j$ in $W^{2,p_0}$ on smaller balls. Thus $Ds_j$ is bounded in
$L^{p_0}$. Failure of Equation~\eqref{reg-drop} gives
$\inf_{B_{1/6}}s_j\to0$. The weak Harnack Inequality with bounded
divergence and scalar errors implies $s_j\to0$ in measure on
$B_{1/3}$; one may use Theorem~5.2 and Corollary~5.4 of
\cite{trudinger1973}, with the radii $1/6$, $1/3$, and $5/6$.

We need more than this convergence in measure. Test
Equation~\eqref{reg-supersolution}, after dropping its nonnegative
Hessian term, with $\eta^2(h-s_j)_+$, where $\eta$ is supported
in $B_{1/3}$. Ellipticity and Young's Inequality give, for fixed
$h>0$,
\begin{equation}
 \int_{\{s_j<h\}}\eta^2|Ds_j|^2
       \le Ch^2\int|D\eta|^2+o_j(1).
 \label{reg-truncated-energy}
\end{equation}
On $\{s_j\ge h\}$, convergence in measure and the uniform
$L^2$ bound for $Ds_j$ give convergence to zero of its $L^1$
norm. On the complementary set, Equation~\eqref{reg-truncated-energy}
and Cauchy--Schwarz bound the $L^1$ norm by $Ch+o_j(1)$.
First let $j\to\infty$, then $h\downarrow0$. We obtain
$Ds_j\to0$ in $L^1(B_{1/4})$. Testing the full
Equation~\eqref{reg-supersolution} with a nonnegative cutoff now
yields
\[
 \int_{B_{1/8}}|\Hess z_j|^2\longrightarrow0.
\]
Since $g_j\to I$ and $Dg_j\to0$, the same is true for $D^2z_j$.
Poincar\'e's Inequality and the uniform Lipschitz bound show, after
passage to a subsequence, uniform convergence on smaller balls to
an affine function. Its slope has length one, because $s_j\to0$
in measure and the gradients converge strongly in $L^2$.
This contradicts failure of Equation~\eqref{reg-affine-alternative}.
At a boundary center the functions are odd, so their affine limit
has zero constant and tangential terms. Increasing $k_*$ if
necessary ensures $k_*>r_*$.
\end{proof}

\begin{lemma}[Improvement of affine approximation]
\label{reg-improvement}
There are $\alpha_2\in(0,\alpha_1)$, $r_2\in(0,1/8)$, and
$\eta_0,h_0>0$ depending only on $c_0$ such that the following
holds. Suppose $g(0)=I$, $|\nabla z|\le8$,
\[
 \|z-L_0\|_{L^\infty(B_1)}\le h\le h_0,
 \qquad \tfrac14\le |DL_0|\le4,
 \qquad
 \|G\|_\infty+\|g-I\|_\infty+\|Dg\|_\infty\le h\eta_0.
\]
Then an affine $L_1$ satisfies
\begin{equation}
 \|z-L_1\|_{L^\infty(B_{r_2})}
       \le h r_2^{1+\alpha_2},\qquad
 |DL_1-DL_0|\le Ch.
 \label{reg-improved-flatness}
\end{equation}
If $z$ and $L_0$ are odd about a reflecting plane through the
center, then $L_1$ can also be chosen odd.
\end{lemma}

\begin{proof}
Set $Y=(z-L_0)/h$ and hold the original coefficients of the
$z$-equation fixed when applying Lemma~\ref{reg-cutoff} to $Y$.
The covariant Hessian of $L_0$ is $-\Gamma DL_0$, so that lemma
gives $\|Y\|_{W^{2,p_0}(B_{7/8})}\le C$. Put
$a=DL_0/|DL_0|$ and $A_0=I-(1-b)a\otimes a$.
The difference between the original axis projector and $a\otimes a$
obeys
\[
 |e\otimes e-a\otimes a|
       \le C\min\{1,h|DY|\}+Ch\eta_0.
\]
This remains valid at a zero gradient, since then $h|DY|=|DL_0|$;
metric changes have been included in the second term. For
$q=2p_0/(p_0-2)>p_0$ and $\vartheta=(p_0-2)/2>0$,
\[
 \|\min\{1,h|DY|\}\|_{L^q(B_{3/4})}
       \le C h^{p_0/q}=Ch^\vartheta.
\]
Consequently H\"older's Inequality, with $1/2=1/q+1/p_0$, gives
\begin{equation}
 A_0^{ij}D_{ij}Y=F_*,\qquad
 \|F_*\|_{L^2(B_{3/4})}\le C(h^\vartheta+\eta_0).
 \label{reg-axis-error}
\end{equation}
Here $F_*$ includes the matrix error times $D^2Y$, the term
$G/h$, and $A^{ij}\Gamma_{ij}^k(D_kL_0/h+D_kY)$.
Thus neither $DY$ nor any derivative of $b$ has been assumed
pointwise bounded.

There is a small zero-Dirichlet correction for
Equation~\eqref{reg-axis-error}. On the convex ball $B_{3/4}$,
\begin{equation}
 \|D^2u\|_2\le\|\Delta u\|_2
       \quad(u\in W^{2,2}\cap W^{1,2}_0).
 \label{reg-miranda-talenti}
\end{equation}
For smooth functions this follows by integrating twice by parts:
$\int((\Delta u)^2-|D^2u|^2)=
\int_{\partial B_{3/4}}H_{\partial B_{3/4}}(\partial_\nu u)^2\ge0$,
with the outward convex convention. Density proves the stated
version. The map
\[
 u\longmapsto\Delta_D^{-1}\bigl(F_*+(1-b)u_{aa}\bigr)
\]
is therefore a contraction, in the Hessian norm, with factor
$1-c_0$. Its fixed point $u$ satisfies $A_0:D^2u=F_*$ and
$\|u\|_{W^{2,2}}\le C\|F_*\|_2$. In dimension three the
embedding $W^{2,2}\hookrightarrow L^\infty$ gives
\[
 \|u\|_\infty\le C(h^\vartheta+\eta_0).
\]
The function $Y_0=Y-u$ solves the fixed-axis equation and has a
uniform $W^{2,2}$ norm. Lemma~\ref{reg-fixed-axis} applies.

Choose $\alpha_2<\alpha_1$, then $r_2$ so small that the Taylor
error $Cr_2^{1+\alpha_1}$ is at most
$\tfrac14r_2^{1+\alpha_2}$. Next choose $\eta_0$, and finally
$h_0$, so that $C(h_0^\vartheta+\eta_0)$ is at most the same
quantity. With
$L_1=L_0+h(Y_0(0)+DY_0(0)\cdot x)$, these bounds prove
Equation~\eqref{reg-improved-flatness}.
If the center is on the boundary plane, $a$ is its normal, $b$ is
even, and $Y$ and $F_*$ are odd. Uniqueness of the correction on
the symmetric ball makes $u$ odd. Hence $Y_0$ is odd and its
Taylor polynomial has only a normal linear term, as required.
\end{proof}

\begin{proof}[Proof of Theorem~\ref{reg-rank-one}]
Choose the constants of Lemma~\ref{reg-improvement} first, and
then decrease $h_*>0$ so that $h_*\le h_0$ and
$Ch_*\sum_{j\ge0}r_2^{j\alpha_2}<1/4$. Use this $h_*$ in
Lemma~\ref{reg-dichotomy}. Set $M=L+Q$; the case $M=0$ is
immediate. Normalize coordinates at each center as above. Choose
a uniform initial radius $\rho_0$ so small that, after dividing
$z$ by $M$ and scaling this radius to one, all errors required by
Lemma~\ref{reg-dichotomy} are at most $\varepsilon_*$, and the
forcing and first-order metric errors are also at most $h_*\eta_0$.

As long as the drop alternative holds, the successive physical
radii and gradient amplitudes are
\begin{equation}
 \rho_j=\rho_0r_*^j,
 \qquad M_j=Mk_*^j,
 \qquad
 z_j(y)=\frac{z(x+\rho_jy)-z(x)}{\rho_j M_j}.
 \label{reg-drop-scales}
\end{equation}
The forcing scales by $\rho_j/M_j$, which decreases because
$r_*<k_*$. The geometric errors decrease with $\rho_j$.
In a reflected patch the plane density is recomputed from
Equation~\eqref{reg-boundary-flux} with the newly normalized
gradient; it is again $O(\rho_j)$. Thus the class is preserved
at every drop, including at centers close to, but not on, the plane.

If drops continue indefinitely, $Dz(x)=0$, and the constant
affine approximations on these balls have errors at most
$C\rho_jM_j$. Otherwise let $m$ be the first index where the
affine alternative holds. At physical radius
$S=\rho_0r_*^{m+1}$, divide by $SM_m$ to obtain affine
approximation with error $h_*$ on the unit ball and slope in
$[1/2,2]$. Apply Lemma~\ref{reg-improvement} repeatedly. At step
$j$ the radius is $r_2^j$ and the flatness is
$h_j=h_*r_2^{j\alpha_2}$. The gradient amplitude $M_m$ is kept
fixed during this iteration; the spatial dilation alone is changed.
The forcing and metric errors decrease as $r_2^j$, so their ratios
to $h_j$ decrease as $r_2^{j(1-\alpha_2)}$. The total change of
the normalized slopes is less than $1/4$, preserving the required
slope range. At boundary centers every affine approximation is odd.

Choose
\[
 0<\alpha\le\min\left\{\alpha_2,
                           \frac{\log k_*}{\log r_*}\right\}.
\]
The drop scales give $M_j\le M(\rho_j/\rho_0)^\alpha$.
After the first affine scale,
\[
 \inf_{L\ \mathrm{affine}}
       \|z-L\|_{L^\infty(B_r(x))}
 \le C M_mS(r/S)^{1+\alpha_2}
 \le CM\rho_0^{-\alpha}r^{1+\alpha}\qquad(r\le S).
\]
Before that scale the constant approximations have the same last
bound. Passing from the discrete radii to arbitrary radii loses
only a fixed factor. Every center therefore has affine
approximations with error $CMr^{1+\alpha}$. Comparing consecutive
approximations on the same ball bounds their slope differences by
$CMr^\alpha$; comparing overlapping balls gives the same bound
between the limiting slopes at two centers. Smoothness identifies
each limiting slope with $Dz$. This proves the first estimate in
Equation~\eqref{reg-main-estimates}, also for the reflected function.

For the second estimate take an affine approximation $L_{x,2r}$
on $B_{2r}(x)$ and apply the exponent-two form of
Lemma~\ref{reg-cutoff} to $z-L_{x,2r}$. Its source is $G$ plus
$A^{ij}\Gamma_{ij}^kD_kL_{x,2r}$, and its slope is bounded by
$CM$. Scaling that estimate back gives
\[
 \|D^2z\|_{L^2(B_r(x))}
 \le C\left(r^{-2}\|z-L_{x,2r}\|_{L^2(B_{2r}(x))}
                      +Mr^{3/2}\right)
 \le CM r^{1/2+\alpha}.
\]
The connection term converting $D^2z$ to $\Hess z$ has squared
integral at most $CM^2r^3$. This proves the claimed Hessian bound.
Restricting the reflected estimate to the original side proves the
boundary assertion.
\end{proof}

\subsection{A bounded scalar equation with natural growth}

\begin{lemma}[H\"older control from a Hessian measure bound]
\label{reg-natural-growth}
Suppose $Z$ is locally Lipschitz, $|Z|\le M$, and satisfies, in
distributions on a ball in $\mathbb R^3$,
\begin{equation}
 \divg(aDZ+B)=e,
 \qquad \lambda I\le a=a^T\le\Lambda I,
 \qquad \|B\|_\infty\le B_0,
 \label{reg-natural-equation}
\end{equation}
where $a$ is measurable and $e$ is a signed Radon measure. Assume
\begin{equation}
 |e|\le C_0(1+|DZ|^2)\,dx+\mu,
 \qquad
 \mu(B_r(x))\le C_\mu r^{1+\beta},
 \qquad 0<\beta\le1,
 \label{reg-natural-measure}
\end{equation}
for all smaller balls, where $\mu$ is nonnegative. Then $Z$ has
a local $C^{0,\theta}$ bound for some $\theta>0$, depending only
on the displayed constants and the patch separation. No bound for
its local Lipschitz constant enters this estimate. The conclusion
also holds at a smooth boundary with bounded total conormal flux,
or with constant Dirichlet value.
\end{lemma}

\begin{proof}
Write $\mathcal L=-\divg(aD)$, choose
$k\ge\max\{1,2C_0/\lambda\}$, and set $V_s=e^{skZ}$ for
$s=\pm1$. The exact distributional chain rule is
\begin{equation}
 \mathcal LV_s
  =\divg(skV_sB)-skV_se
    -k^2V_s\bigl(aDZ\cdot DZ+B\cdot DZ\bigr).
 \label{reg-exponential-identity}
\end{equation}
It is valid for the present weak equation: the exponential is
locally Lipschitz and can be approximated uniformly and in
$W^{1,2}$ in the test, so the measure term converges as well.
The inequality
$|B\cdot DZ|\le(\lambda/2)|DZ|^2+B_0^2/(2\lambda)$ and
Equation~\eqref{reg-natural-measure} imply
\begin{equation}
 \mathcal LV_s\le\divg F_s+\nu,
 \qquad F_s=skV_sB,
 \qquad
 \nu=e^{kM}\left(kC_0+\frac{k^2B_0^2}{2\lambda}\right)dx
                     +ke^{kM}\mu.
 \label{reg-exponential-subsolution}
\end{equation}
In particular $F_s$ is bounded and $\nu$ is a positive measure
with the same admissible growth, in addition to a bounded density.

\emph{The correction estimate.}
We record the small correction estimate on $B_R$:
\begin{equation}
 \mathcal Lw_s=\divg F_s+\nu,
 \qquad w_s\in W^{1,2}_0(B_R),
 \qquad \|w_s\|_\infty\le\varepsilon_R=CR^\beta.
 \label{reg-small-correction}
\end{equation}
For the bounded vector term, testing the Dirichlet equation with
$(w-t)_+$ gives, for $p>3$,
\[
 \int|D(w-t)_+|^2
       \le C\|F_s\|_p^2|\{w>t\}|^{1-2/p}.
\]
Sobolev's Inequality then yields
\[
 |\{w>h\}|
 \le C\|F_s\|_p^6(h-t)^{-6}
                 |\{w>t\}|^{3-6/p},\qquad h>t.
\]
The exponent on the right is greater than one. Iteration, also
for $-w$, gives $\|w\|_\infty\le
CR^{1-3/p}\|F_s\|_p\le CR$.

For the measure term the positive scalar Dirichlet Green function
satisfies $G_{B_R}(x,y)\le C|x-y|^{-1}$. This is the scalar
measurable-coefficient Green bound: extend $a$ elliptically to
$\mathbb R^3$, use the fundamental-solution bound in
Theorem~3.1 of \cite{hofmannkim2007}, and compare the Dirichlet
Green function by the weak maximum principle. The scalar
De Giorgi--Nash estimate supplies the local H\"older hypothesis
of that theorem. Dyadic annuli and
Equation~\eqref{reg-natural-measure} give
\[
 \sup_x\int_{B_R}G_{B_R}(x,y)\,d\nu(y)
                  \le C(R^2+R^\beta).
\]
This also constructs the needed weak correction without assuming
measure solvability. Approximate the restricted positive measure
by smooth positive measures, preserving the ball-growth bound.
The resulting potentials have the displayed uniform supremum
bound, and their energy is bounded by their supremum times the
total mass. A weak $W^{1,2}_0$ limit solves the measure equation
and retains the bound. Adding the vector and measure corrections
proves Equation~\eqref{reg-small-correction}.

\emph{Oscillation decay.}
Let $M_R=\sup_{B_R}Z$, $m_R=\inf_{B_R}Z$, and set
\begin{equation}
 H_+=e^{kM_R}-e^{kZ}+w_++\varepsilon_R,
 \qquad
 H_-=e^{-km_R}-e^{-kZ}+w_-+\varepsilon_R.
 \label{reg-deficits}
\end{equation}
These functions are nonnegative weak supersolutions of
$\mathcal LH_\pm\ge0$. Each exponential deficit is bounded
above and below by positive constants, depending on $kM$, times
$M_R-Z$ or $Z-m_R$, respectively. One of the two midpoint sets
\[
 \{Z\le(M_R+m_R)/2\}\cap B_{R/3},\qquad
 \{Z\ge(M_R+m_R)/2\}\cap B_{R/3}
\]
has at least half the volume of $B_{R/3}$. Apply weak Harnack to
the corresponding $H_\pm$. In Theorem~5.2 of
\cite{trudinger1973}, take the inner radius $R/6$: its averaging
ball is $B_{R/3}$ and its required nonnegativity ball is
$B_{5R/6}$. Thus the infimum on $B_{R/6}$ is at least
$c(M_R-m_R)$. Removing the correction in
Equation~\eqref{reg-deficits} costs at most $2\varepsilon_R$.
It follows that either the upper endpoint decreases or the lower
endpoint increases by $c(M_R-m_R)-C\varepsilon_R$, whence
\begin{equation}
 \operatorname{osc}_{B_{R/6}}Z
       \le(1-c)\operatorname{osc}_{B_R}Z+CR^\beta.
 \label{reg-oscillation-decay}
\end{equation}
Iteration gives any sufficiently small positive exponent below
both $\beta$ and $-\log(1-c)/\log6$.

\emph{Boundary reflection.} Flatten a boundary face and put
$S=\operatorname{diag}(1,1,-1)$. For an even extension of $Z$,
extend $a^-=Sa^+S$ and $B^-=SB^+$. If $q$ is the third
component of the original total flux on $x_3=0$, the reflected
equation has the additional density $2q\,\delta_{\{x_3=0\}}$.
It is bounded by the prescribed total conormal bound, after the
coordinate volume factors are included. For a constant Dirichlet
value, extend $Z$ minus that value oddly and use $a^-=Sa^+S$,
$B^-=-SB^+$ and the odd bulk source. The normal flux is then
continuous, with no plane term. The reflected matrix remains
measurable and uniformly elliptic. A bounded plane density has
ball mass $O(r^2)$, which is allowed in
Equation~\eqref{reg-natural-measure} because $\beta\le1$.
This proves the boundary assertion.
\end{proof}

\subsection{Classical bootstrap from bounded scalar variables}

\begin{proposition}[The regularity interface for the coupled system]
\label{reg-coupled-bootstrap}
Consider a family of smooth solutions $(f,Z)$ on smooth finite
domains with uniformly controlled interior and boundary patches.
Assume $|f|+|Df|+|Z|\le M$ and the following structural bounds,
uniformly in the family:
\begin{enumerate}[label=(\roman*)]
\item The $f$-equation has the form of
Equation~\eqref{reg-trace-equation}, with $b\ge c_0>0$ and
bounded $G$. Its actual matrix and source are smooth functions of
$(x,Z,f,Df)$ on the bounded variable range, including at $Df=0$.
The value of $f$ on each boundary face is constant.
\item After multiplication by the coordinate volume density, the
$Z$-equation is
\begin{equation}
 \partial_i\bigl(a^{ij}(x,Z,f,Df)\partial_jZ
                           +B^i(x,Z,f,Df)\bigr)=Q,
 \qquad
 |Q|\le C(1+|DZ|^2+|D^2f|^2),
 \label{reg-coupled-Z}
\end{equation}
where $a$ is symmetric and uniformly elliptic and $B$ is bounded.
The functions $a,B$ and $Q$, as a function also of $(DZ,D^2f)$,
are smooth on the ranges on which they are used.
\item Each boundary component for $Z$ has either constant
Dirichlet data, or bounded total conormal flux. In the latter case
its equation can also be solved as
\begin{equation}
 \partial_\nu Z=H(x,Z,f,Df),
 \label{reg-normal-data}
\end{equation}
where $H$ is smooth on the bounded variable range. These faces
are disjoint and have no edges where the type of data changes.
\end{enumerate}
Assume that the derivatives of the displayed smooth functions and
the geometry are bounded to the order needed for each stated
estimate. Then $f$ and $Z$ have uniform local classical estimates
of every such finite order, including at the boundary. The constants
are independent of the size of a finite exhaustion domain if the
patch geometry and all the displayed data bounds are independent
of that size.
\end{proposition}

\begin{proof}
Theorem~\ref{reg-rank-one} gives a uniform $C^{1,\alpha}$ bound
for $f$ and
\begin{equation}
 \int_{B_r(x)}|D^2f|^2\le Cr^{1+2\alpha}.
 \label{reg-bootstrap-measure}
\end{equation}
Coordinate connection terms do not affect this estimate. Apply
Lemma~\ref{reg-natural-growth} to Equation~\eqref{reg-coupled-Z}
with $\mu=C|D^2f|^2dx$ and
$\beta=\min\{2\alpha,1\}$. At a flux boundary include the
bounded reflected plane density. We obtain a uniform H\"older
bound for $Z$. Smooth composition in the bounded variables now
makes the leading matrix and source of the $f$-equation H\"older.
The interior and constant-Dirichlet Schauder estimates give
$f\in C^{2,\theta}$ for some uniform $\theta>0$.

Expand Equation~\eqref{reg-coupled-Z} in nondivergence form.
Its leading matrix is uniformly elliptic and $C^{0,\theta}$.
Differentiating its coefficients produces terms at most quadratic
in $DZ$, because $D^2f$ is now bounded. Thus its right-hand side
is bounded by $C(1+|DZ|^2)$. Fix $p>3$. Linear local $W^{2,p}$
estimates, with either constant Dirichlet or the normal data in
Equation~\eqref{reg-normal-data}, give on nested patches
$U_x\Subset V_x\Subset V_x^+$
\begin{equation}
 \|Z\|_{W^{2,p}(U_x)}
 \le C\left(1+\|DZ\|_{L^{2p}(V_x)}^2
                         +\|Z\|_{W^{1,p}(V_x^+)}\right).
 \label{reg-linear-W2p}
\end{equation}
The patch size here is chosen from the already established
H\"older modulus of the leading matrix. The normal boundary
operator is uniformly oblique; the local a priori estimate is,
for example, a smooth-boundary specialization of Theorem~2.3
of \cite{dongli2022}, followed by localization. Only its estimate
is used, not solvability for a pure oblique problem. The boundary
norm is controlled by
\begin{equation}
 \|H(x,Z,f,Df)\|_{W^{1-1/p,p}(\partial V_x)}
       \le C\left(1+\|Z\|_{W^{1,p}(V_x^+)}\right).
 \label{reg-boundary-trace}
\end{equation}
Indeed, extend the smooth expression for $H$ through the collar,
differentiate it there, use the already bounded first two
derivatives of $f$, and apply the trace theorem. This explains
the last term in Equation~\eqref{reg-linear-W2p} without assuming
a gradient bound for $Z$.

We give the absorption of the quadratic gradient term. On balls
or half-balls of radius $h$, the scaled interpolation estimate is
\begin{equation}
 \|DZ\|_{L^{2p}(B_h)}^2
 \le C\operatorname{osc}_{B_{2h}}Z\,
                 \|D^2Z\|_{L^p(B_{2h})}
       +Ch^{3/p-2}(\operatorname{osc}_{B_{2h}}Z)^2.
 \label{reg-interpolation}
\end{equation}
One way to obtain it is to subtract a constant, integrate
$\int\eta^{2p}|DZ|^{2p}$ by parts against $Z$, and use
H\"older and Young Inequalities on the terms containing $D^2Z$
and $D\eta$. At the boundary use a bounded extension operator
on a half-ball, bounded also in $L^\infty$; its lower derivative
terms are absorbed by the same interpolation. Scaling gives
exactly the second power of $h$ displayed in
Equation~\eqref{reg-interpolation}.

Write $[Z]_{C^{0,\theta}}\le H_0$. Cover $V_x$ by radius-$h$
patches of bounded overlap, with their doubles in $V_x^+$.
Take the $\ell^p$ sum of Equation~\eqref{reg-interpolation}.
There are $O(h^{-3})$ patches, so
\begin{equation}
 \|DZ\|_{L^{2p}(V_x)}^2
 \le CH_0h^\theta\|D^2Z\|_{L^p(V_x^+)}
                         +CH_0^2h^{2\theta-2}.
 \label{reg-gradient-absorption}
\end{equation}
Also, for every $\eta>0$,
$\|Z\|_{W^{1,p}(V_x^+)}
\le\eta\|D^2Z\|_{L^p(V_x^{++})}+C_\eta$
on a further uniformly enlarged patch. Such enlarged patches
are covered by a bounded number of the $U_y$. On a fixed
smooth finite-domain solution let
$T=\sup_y\|Z\|_{W^{2,p}(U_y)}$, which is finite before making
any uniform assertion. Equations~\eqref{reg-linear-W2p} and
\eqref{reg-gradient-absorption} imply
\begin{equation}
 T\le C(H_0h^\theta+\eta)T+C(h,\eta).
 \label{reg-global-absorption}
\end{equation}
The radius for the linear estimate was fixed first. Now choose
the auxiliary covering radius $h$, and then $\eta$, so that the
coefficient of $T$ is less than $1/2$. This proves a uniform
$W^{2,p}$ bound, and hence $Z\in C^{1,1-3/p}$.

The source in the nondivergence $Z$-equation is now H\"older;
the normal data have the corresponding first derivative
regularity. Schauder estimates give $Z\in C^{2,\theta'}$.
The $f$-equation then gains another derivative. Repeating this
argument gives every finite order supported by the data.
All coverings and constants were uniform, so this conclusion
has the asserted independence from the exhaustion radius.
\end{proof}

\begin{remark}
The hypotheses in Proposition~\ref{reg-coupled-bootstrap} are
an interface to be checked from the actual deformation equations.
In particular the lower bound for $b$, the bounded variables,
the quadratic source estimate, and the solvable normal boundary
condition are required before applying the proposition. The
argument supplies classical bounds on fixed-size patches; it
does not by itself give decay at infinity. Uniform estimates of
arbitrarily high order on an end require correspondingly high
derivative bounds for the reduced background data.
\end{remark}
